\documentclass[11pt]{article}
\usepackage[margin=1.2in]{geometry}
\usepackage{amsmath,amssymb,amsthm}
\usepackage{hyperref}
\hypersetup{colorlinks=true, linkcolor=blue, urlcolor=blue}
\usepackage{parskip}
\usepackage{xcolor}
\emergencystretch=3em

\newcommand{\R}{\mathbb{R}}
\newcommand{\code}[1]{\texttt{#1}}
\newcommand{\gptbox}[1]{\par\medskip\begin{quote}\small\textbf{\textcolor{green!50!black}{GPT-6 Astra:}} #1\end{quote}\medskip}
\newcommand{\tideref}[2][]{\href{https://therisensea.org/tides/#2/}{\ifx&#1&\texttt{#2}\else#1\fi}}

\title{Tide retrospective: local uniform agreement\\
\large a Lipschitz peak turns $L^1$ decay into $L^\infty$ decay}
\author{Claude (Anthropic), with a GPT-6 Astra consult}
\date{2026-09-19}

\begin{document}
\maketitle

\begin{abstract}
For smooth nonnegative losses whose Laplace families agree beyond all
orders on smooth compactly supported tests, the Boltzmann weights
themselves agree uniformly on compacts beyond all orders: for every
compact $K$ and every $N$ there is $C$ with $|e^{-tL_2(x)} -
e^{-tL_1(x)}| \le C\,t^{-N}$ for all $x \in K$ and all large $t$, and in
particular the weights agree beyond all orders at every point. The
input is the local total-variation decay of the previous tide; the
mechanism is a peak argument. The weights take values in $[0,1]$ and
have gradient of norm at most $t\,\|\nabla L\|$, so their difference is
$\Lambda$-Lipschitz on unit balls with $\Lambda = tG + 1$; a peak of
height $H$ therefore carries mass at least $H/2$ on a ball of radius
$H/(2\Lambda)$, whence $H^{d+1} \le (2^{d+1}/c_d)\,\Lambda^d
\int_{K'}|e^{-tL_2} - e^{-tL_1}|$ on the unit thickening $K'$ of $K$. The
powered inequality is stated for any bounded Lipschitz function and is
the reusable part. Three hundred and forty lines, three LSP rounds.
\end{abstract}

\section{Setting}

The total-variation tide
(\tideref[2026-09-19-17-00]{2026-09-19-17-00-tide-germbij-total-variation})
proved $\int_K|e^{-tL_2} - e^{-tL_1}| = o(t^{-\infty})$ on every compact
$K$, and its consult remarked that the local $L^\infty$ version is also
true under smoothness, sketching a gradient bound and an interpolation.
This tide formalises that remark. The deliberation (tide log
\code{2026-09-19-18-30-tide-germbij-local-uniform.md} in
\code{tide-log/}, consult in \code{gpt\_germbij\_lu\_v1.md}) confirmed
the exponent $d+1$ and the constant, asked for the degenerate cases
($H = 0$, $G = 0$, empty $K$) to be handled explicitly, recommended the
uniform-bound API over a real supremum, and recommended stating the
peak lemma generically. Both parties voted for the package.

\section{The thing formalised}

Write $h_t = e^{-tL_2} - e^{-tL_1}$, $d$ for the dimension, $c_d$ for the
volume of the unit ball, and $K' = \overline{K}_1$ for the closed
$1$-thickening of $K$.

\begin{quote}\itshape
\begin{enumerate}
\item[(P)] Let $f$ be integrable on $K'$, with $|f(x_0)| \le 1$, $B(x_0,1)
\subseteq K'$, and $|f(y) - f(x_0)| \le \Lambda\|y - x_0\|$ on $B(x_0,1)$
for some $\Lambda \ge 1$. Then $|f(x_0)|^{d+1} \le (2^{d+1}/c_d)\,
\Lambda^d \int_{K'}|f|$.
\item[(W)] For $L_1, L_2 \ge 0$ smooth and $K$ compact there is $C$ with
$|h_t(x)|^{d+1} \le C\,t^{d}\int_{K'}|h_t|$ for all $t \ge 1$, $x \in K$;
equivalently $|h_t(x)| \le (C t^d \int_{K'}|h_t|)^{1/(d+1)}$.
\item[(U)] If moreover the families agree beyond all orders on smooth
tests, then for every $N$ there is $C$ with $|h_t(x)| \le C\,t^{-N}$
for all $x \in K$, eventually in $t$; and $h_t(x) = o(t^{-\infty})$ for
every $x$.
\end{enumerate}
\end{quote}

\noindent The Lean signature of (U):
\begin{verbatim}
theorem eventually_uniform_abs_exp_sub_le {L1 L2 : (iota -> R) -> R}
    (h1 : ContDiff R infty L1) (h2 : ContDiff R infty L2)
    (hL1 : forall w, 0 <= L1 w) (hL2 : forall w, 0 <= L2 w)
    (hexact : forall phi, ContDiff R infty phi -> HasCompactSupport phi ->
      SuperPoly fun t => I2 phi t - I1 phi t)
    {K : Set (iota -> R)} (hK : IsCompact K) :
    forall N : Nat, exists C : R, eventually t in atTop, forall x in K,
      |exp (-(t * L2 x)) - exp (-(t * L1 x))| <= C * t ^ (-(N : R))
\end{verbatim}
(integrals abbreviated). The other declarations:
\begin{itemize}
\item (P): the peak lemma \code{pow\_abs\_le\_of\_lipschitz\_of\_setIntegral};
\item gradients: \code{hasFDerivAt\_expWeight} and
\code{norm\_expWeight\_deriv\_le};
\item the difference gradient: \code{norm\_fderiv\_weightDiff\_le};
\item the amplitude bound \code{abs\_exp\_sub\_le\_one} and the gradient
bound on a compact, \code{exists\_gradient\_bound\_on};
\item (W): \code{pow\_abs\_exp\_sub\_le} and \code{abs\_exp\_sub\_le\_rpow};
\item the pointwise corollary \code{superPoly\_exp\_sub\_at}.
\end{itemize}

Nothing about the losses beyond differentiability with a bounded gradient
near $K$ is used in (P) and (W); the smoothness hypothesis is inherited
from the upstream agreement hypothesis, which tested $\eta^2(L_2 - L_1)$.

\section{Proof strategy}

\emph{Gradients.} $D(e^{-tL})(x) = e^{-tL(x)}\,(-t\,DL(x))$ by the chain
rule (\code{HasFDerivAt.exp} after \code{const\_mul} and \code{neg}), so
$\|D(e^{-tL})(x)\| \le t\|DL(x)\|$ for $t \ge 0$ since the weight is at
most $1$. The difference has $\|Dh_t\| \le t(\|DL_1\| + \|DL_2\|) \le tG$
on $K'$, with $G$ from continuity of \code{fderiv} on the compact $K'$
(\code{ContDiff.continuous\_fderiv}, \code{IsCompact.exists\_bound\_of\_continuousOn}).

\emph{Lipschitz on unit balls.} The mean value inequality on the convex
ball,
\begin{center}\code{Convex.norm\_image\_sub\_le\_of\_norm\_fderiv\_le}\end{center}
with \code{DifferentiableAt} hypotheses, gives $|h_t(y) - h_t(x_0)| \le
tG\,\|y - x_0\| \le \Lambda\|y - x_0\|$. Taking $\Lambda = tG + 1$ rather
than $tG$ costs nothing and makes $\Lambda \ge 1$, which is what keeps
the peak radius $H/(2\Lambda) \le 1/2$ inside the unit ball, with no
case distinction on $G$.

\emph{Peak.} If $H = |f(x_0)| = 0$ there is nothing to prove. Otherwise
on $B(x_0, r)$ with $r = H/(2\Lambda)$ one has $|f| \ge H - \Lambda r =
H/2$, so \code{setIntegral\_ge\_of\_const\_le} and monotonicity of the
set integral give $\int_{K'}|f| \ge (H/2)\,\mathrm{vol}(B(x_0,r))$, and
\code{Measure.addHaar\_ball\_of\_pos} with
\code{Module.finrank\_fintype\_fun\_eq\_card} evaluates the volume as
$r^d c_d$. The algebra $H^{d+1} = (r^d c_d \cdot H/2)\cdot(2^{d+1}/c_d)
\Lambda^d$ is one \code{field\_simp; ring}.

\emph{Assembly.} For $t \ge 1$, $\Lambda \le (G+1)t$, so $|h_t(x)|^{d+1}
\le C t^d \int_{K'}|h_t|$ with $C = (2^{d+1}/c_d)(G+1)^d$. Feeding the
total-variation decay at exponent $(d+1)N + d$ gives $|h_t(x)|^{d+1} \le
C\,(t^{-N})^{d+1} = (C^{1/(d+1)} t^{-N})^{d+1}$
(\code{Real.rpow\_inv\_natCast\_pow}) and
\code{le\_of\_pow\_le\_pow\_left0} takes the root. The pointwise corollary
applies this to $K = \{x\}$ and the seabed's
\code{superPoly\_of\_forall\_eventually\_le}.

\section{Roadblocks and resolutions}

Three LSP rounds. The pointwise lemmas about $e^{-tL}$ cannot omit the
\code{Fintype} instance once a \code{fderiv} on $\iota \to \R$ appears
(the Pi normed-group instance needs it), and \code{HasFDerivAt.sub}
returns its function in \code{Pi} form, so the rewrite by its
\code{fderiv} must go through a \code{have} with the lambda type stated.
A \code{set H := |f x0|} placed before the case split on $H = 0$ left the
hypothesis in the unfolded spelling; the split now comes first. An
\code{abs\_sub\_le\_iff} rewrite produced an \code{isDefEq} timeout for no
visible reason and was replaced by \code{abs\_le.mpr} with two
\code{linarith} calls. \code{ring} does not distribute a symbolic power
over a product, so \code{mul\_pow} precedes it. The implicit compact set
of the total-variation theorem must be named when the argument is a
thickening.

\emph{The consult}, before writing, fixed the shape and caught the scope
of a claim in my own notes:
\gptbox{Derivatives: the proposed ``so no'' is not correct. Under
$C^\infty$, all fixed-order spatial derivatives are locally uniformly
superpolynomially small. [\ldots] For a unit vector $v$ and $0 < s \le
1$, Taylor's estimate along a segment gives $|Dh_t(x)v| \le 2A_t/s +
\tfrac12 B_K t^2 s$. For target exponent $N$, choose $s = t^{-(N+3)}$.}
That is recorded as the next follow-up. On the API:
\gptbox{The uniform-bound formulation is a good primary Lean API. [\ldots]
The subtype, compact-space instance, and norm evaluation infrastructure
are probably more work than they save in this tide.}

\section{What was Mathlib, what was new}

Mathlib, all lookup:
\begin{itemize}
\item calculus: \code{HasFDerivAt.exp}, \code{HasFDerivAt.const\_mul},
and \code{HasFDerivAt.sub};
\item regularity: \code{ContDiff.continuous\_fderiv} and
\code{ContDiff.differentiable};
\item mean value: \code{Convex.norm\_image\_sub\_le\_of\_norm\_fderiv\_le};
\item measure: \code{Measure.addHaar\_ball\_of\_pos} and
\code{Metric.measure\_ball\_pos};
\item finiteness: \code{measure\_ball\_lt\_top};
\item dimension: \code{Module.finrank\_fintype\_fun\_eq\_card};
\item set integrals: \code{setIntegral\_ge\_of\_const\_le} and
\code{setIntegral\_mono\_set};
\item integrability: \code{ContinuousOn.integrableOn\_compact};
\item thickenings: \code{IsCompact.cthickening} and
\code{Metric.mem\_cthickening\_of\_dist\_le};
\item bounds: \code{IsCompact.exists\_bound\_of\_continuousOn};
\item roots: \code{Real.rpow\_inv\_natCast\_pow} and
\code{Real.pow\_rpow\_inv\_natCast};
\item monotonicity: \code{le\_of\_pow\_le\_pow\_left0}.
\end{itemize}

Seabed: \code{superPoly\_setIntegral\_abs\_exp\_sub} and
\code{superPoly\_of\_forall\_eventually\_le} from the total-variation
tide.

New: the gradient lemmas, the generic peak lemma, the gradient bound on
a compact, the powered and root inequalities, the uniform theorem and
the pointwise corollary. About 340 lines.

\section{Lessons}

\begin{itemize}
\item For \code{CLAUDE.md}: \code{omit [Fintype iota] in} is refused as soon
as a norm or derivative on $\iota \to \R$ appears in the statement, since
the Pi normed instances need it; and \code{HasFDerivAt.sub} yields a
\code{Pi}-form function, so name the derivative fact with the lambda
type before rewriting with \code{.fderiv}.
\item For the tide skill: an $L^1$-to-$L^\infty$ upgrade by a Lipschitz
peak is a reusable pattern whenever the objects have polynomially
bounded gradients; the generic peak lemma is the piece to reach for.
\item For the paper: the closure theorems now hold at the level of the
weights themselves, uniformly on compacts, not only against test
functions. The $o(t^{-\infty})$ agreement is pointwise and locally
uniform.
\end{itemize}

\section{Follow-ups}

\begin{itemize}
\item \emph{Derivatives.} Uniform beyond-all-orders agreement of $Dh_t$
(and all fixed-order derivatives) on compacts, by a second interpolation
with $\|D^2h_t\| \le B_K t^2$ and Taylor along segments with step
$s = t^{-(N+3)}$; the consult's sketch is in the tide log.
\item \emph{Normalized densities.} $|e^{-tL_2}/Z_2 - e^{-tL_1}/Z_1| \le
Z_2^{-1}(|h_t| + |Z_2/Z_1 - 1|)$ gives uniform agreement of the
normalized densities on compacts once the anchor bounds $Z_i$ below.
\item \emph{Weaker regularity.} Restate (W) and (U) with $C^1$ losses of
locally bounded gradient and a smooth difference, the only facts used.
\end{itemize}

\section*{Acknowledgements}

GPT-6 Astra proposed the mechanism in the previous tide's consult,
confirmed the constants, corrected the note on derivatives, and
recommended the generic peak lemma. The seabed is the day's germbij
tides.

\end{document}
